% Source: https://artofproblemsolving.com/community/c3h1501818_similar_triangle
% Stable identity: aops-topic-1501818-q01
Let $X$ be the intersection of $CC_1$ and $AA_1$, where $C_1$ is the midpoint of $AB$:

\begin{asy}
usepackage("amsmath");
size(5.5cm);
defaultpen(fontsize(10pt));
pair A=(1,3),B=(0,0),C=(4,0);
pair A1=B+2/3*(C-B), C1=(A+B)/2, X=A+3/4*(A1-A);
draw(A--B--C--cycle,linewidth(1));
draw(A--A1,blue+linewidth(0.9)); draw(C--C1,red+linewidth(0.9));
pair t1=(A+C1)/2, t2=(C1+B)/2, u=0.09*unit(rotate(90)*(A-B));
draw(t1-u--t1+u); draw(t2-u--t2+u);
dot(A);dot(B);dot(C);dot(A1);dot(C1);dot(X);
label("$A$",A,N);label("$B$",B,SW);label("$C$",C,SE);label("$A_1$",A1,S);label("$C_1$",C1,W);label("$X$",X,dir(80));
label("$2$",(B+A1)/2,S,fontsize(9pt)); label("$1$",(A1+C)/2,S,fontsize(9pt));
label("\mbox{}",(min(currentpicture,user=true)+max(currentpicture,user=true))/2,black); // keeps the text after the figure black when LaTeX typesets the labels inline
\end{asy}

Apply Menelaus' theorem to triangle $ABA_1$ and the line through $C_1$, $X$, $C$ (with $C$ on line $BA_1$ extended):
\[
\dfrac{AC_1}{C_1B}\cdot\dfrac{BC}{CA_1}\cdot\dfrac{A_1X}{XA}=1 .
\]
Here $\dfrac{AC_1}{C_1B}=1$ and $\dfrac{BC}{CA_1}=\dfrac{3}{1}$, so $\dfrac{A_1X}{XA}=\dfrac{1}{3}$.

\emph{Check with mass points.} Put mass $1$ at $B$ and $2$ at $C$, so that $A_1$ (with $BA_1:A_1C=2:1$) is their balance point, with mass $3$. Put mass $1$ at $A$, so that $C_1$ balances $A$ and $B$. Then $X$ is the balance point of $A$ (mass $1$) and $A_1$ (mass $3$), so $AX:XA_1=3:1$.

The median divides $AA_1$ in the ratio $\boxed{AX:XA_1=3:1}$, counting from $A$. (It also bisects itself: $CX=XC_1$.)
